%! AP Statistics — Unit 1, Lesson 1.4 — Student Packet Cover Sheet
\documentclass[10pt]{article}
\usepackage{apstats-article}
\usepackage{apstats-boxes}

%=============================================================================
\begin{document}

% ── Full-bleed navy banner ────────────────────────────────────────────────────
\begin{tikzpicture}[remember picture, overlay]
  \fill[navy] (current page.north west)
    rectangle ([yshift=-0.9in]current page.north east);
\end{tikzpicture}
\vspace{-0.8in}

\begin{center}
  {\color{white}\textbf{\LARGE AP Statistics}}\\[4pt]
  {\color{skymid}\large\textbf{Unit 1: Exploring One-Variable Data}}\\[6pt]
  {\color{white}\textbf{\large Lesson 1.4 \quad Representing a Categorical Variable with Graphs}}
\end{center}

\vspace{0.22in}

\namedateperiod

\vspace{0.18in}

% ── LEARNING TARGETS ─────────────────────────────────────────────────────────
\begin{learningtargetbox}
\small
\begin{enumerate}[leftmargin=1.5em, itemsep=5pt]
  \item \ldots{} construct a correctly labeled \textit{bar chart} (both frequency
        and relative frequency versions) from a categorical data table.
  \item \ldots{} construct a \textit{pie chart} by computing each sector's degree
        measure as (relative frequency) $\times\,360^\circ$.
  \item \ldots{} interpret bar charts and pie charts in context, identifying the
        most and least common categories and describing the overall distribution.
  \item \ldots{} identify misleading features in a graph (e.g., non-zero y-axis,
        3-D distortion, inconsistent scale) and explain how to fix them.
\end{enumerate}
\end{learningtargetbox}

\vspace{0.14in}

% ── TABLE OF CONTENTS ────────────────────────────────────────────────────────
\begin{tocbox}
\small
\renewcommand{\arraystretch}{1.5}
\begin{tabularx}{\linewidth}{@{} c l X r @{}}
\rowcolor{sky}
\textbf{\#} & \textbf{Component} & \textbf{Description} & \textbf{Score} \\
\midrule
1 & Warm-Up        & Spiral review + sketch a visual display                  & \blank{1.2cm} \\
2 & Guided Notes   & Bar charts, pie charts, misleading graphs                 & \blank{1.2cm} \\
3 & Group Activity & Construct and interpret a relative frequency bar chart    & \blank{1.2cm} \\
4 & Exit Ticket    & Sketch a bar chart + describe the distribution            & \blank{1.2cm} \\
5 & Homework       & Chart construction, misleading graph critique, AP-style   & \blank{1.2cm} \\
\midrule
  & \multicolumn{2}{r}{\textbf{Total}} & \blank{1.2cm} \\
\end{tabularx}
\end{tocbox}

\vspace{0.14in}

% ── KEEP IN MIND ─────────────────────────────────────────────────────────────
\begin{remindbox}
\small
Bar charts and pie charts display the \emph{same} categorical distributions
you summarized with tables in Lesson 1.3 --- just in visual form. A graph that
looks impressive can still be dishonest; always check the axis.

\vspace{6pt}
\begin{tabularx}{\linewidth}{@{} >{\centering\arraybackslash\columncolor{goldbg}}p{0.06\linewidth}
                                    X @{}}
\textbf{\large!} &
\textbf{A bar chart's y-axis must start at 0.}\enspace
A non-zero axis exaggerates differences between bars and can make small
differences look dramatic. \\[6pt]
\textbf{\large!} &
\textbf{Bars in a bar chart do \emph{not} touch.}\enspace
Categories are separate; gaps signal that the variable is categorical,
not continuous (unlike histograms in Lesson~1.5). \\[6pt]
\textbf{\large!} &
\textbf{Every graph needs a title and labeled axes.}\enspace
On the AP exam, an unlabeled graph earns partial or no credit even if
the bars are the correct heights. \\[6pt]
\textbf{\large!} &
\textbf{Use relative frequency to compare; use counts when the total matters.}\enspace
Bar charts with relative frequency on the y-axis are preferred for
comparing groups of different sizes. \\
\end{tabularx}
\end{remindbox}

\end{document}
